\section{Results}
Our experiments on the Turtle Soup QA task involve two main approaches: few-shot prompting and fine-tuning. We evaluated both open-source Large Language Models (LLMs) and proprietary models from OpenAI. The results reveal interesting patterns in model performance and the effects of different training strategies.

\subsection{Few-shot Prompting Performance}
Few-shot prompting was applied to both open-source LLMs and OpenAI models. The results, as shown in Table~\ref{tab:performance}, reveal several key insights:

\subsubsection{Competitive Performance of Open-Source Models}
Contrary to initial expectations, open-source models demonstrated performance levels competitive with those of OpenAI models. Notably, Qwen2.5-72B achieved an F1 score of 80.88\% with 10-shot prompting, surpassing GPT-4o's performance of 80.36\% with 10-shot prompting. This suggests that the gap between proprietary and open-source models is narrowing in complex reasoning tasks.

\subsubsection{Unexpected Performance of o1-preview}
Interestingly, the o1-preview model, which is purported to have superior reasoning capabilities, exhibited lower F1 scores compared to GPT-4o. The highest F1 score for o1-preview was 77.08\% at 50-shot, while GPT-4o achieved 80.36\% at 10-shot. This unexpected result calls for further investigation into the specific strengths and weaknesses of these models, particularly in lateral thinking tasks.

% \subsubsection{Model Size and Performance Correlation}
% Generally, larger models exhibited superior performance. For example, Llama3.1-70B outperformed its 8B counterpart in most shot settings, with a peak F1 score of 75.67\% at 10-shot compared to the 8B model's peak of 76.97\% at 30-shot. This trend is consistent with the common observation that increased model size often correlates with enhanced performance in complex tasks.

\subsubsection{Missing Data Points}
It is important to note that our results include some missing data points, particularly for larger models and higher shot settings. These gaps are primarily due to two key factors:

\begin{enumerate}
    \item \textbf{GPU Memory Limitations:} We encountered CUDA out-of-memory errors with some of the larger models at higher shot settings. For instance, we were unable to obtain data for Llama3.1-70B beyond the 30-shot setting. These limitations restricted our ability to evaluate these models at higher shot counts.

    \item \textbf{Time Constraints:} The evaluation process for some models at higher shot counts was extremely time-consuming. In certain cases, the allocated job time expired before the evaluation could be completed. This issue is evident with models like InternLM2.5-20B, for which we only have data for the 0-shot setting. Similarly, we lack data for Qwen2-72B and Qwen2.5-72B beyond the 20-shot setting, and for Mistral-7B at the 50-shot setting, due to the same reason.
\end{enumerate}

These findings underscore the significant computational resources and time required for a comprehensive assessment of the largest language models, particularly in few-shot scenarios.

\subsubsection{Performance Trends Across Shot Counts}
Despite some missing data points, several trends in model performance across different shot counts can be observed:

\begin{itemize}
    \item \textbf{Variability in Optimal Shot Count:} Different models appear to reach their peak performance at varying shot counts. For instance, GPT-4o achieved its highest F1 score of 80.36\% at 10-shot, whereas o1-preview peaked at 77.08\% at 50-shot. This suggests that the optimal number of examples may vary depending on the model's architecture and size.
    
    \item \textbf{Inconsistent VCR:} The Valid Classification Ratio (VCR) demonstrates significant variability across models and shot counts, particularly among open-source models. For instance, Mistral-7B-v0.3 shows extremely low VCR values (ranging from 1.17\% to 14.20\%) compared to other models. This inconsistency in VCR highlights a critical challenge in few-shot learning scenarios and underscores the need for exploring fine-tuning approaches to improve model performance and stability.
\end{itemize}

\begin{table*}[htbp]
\centering
\small
\caption{Performance of Language Models on the Turtle Soup QA Task Across Different Shot Counts (F1 Score and Valid Classification Ratio in \%). \small{The top metrics for each model across different shot counts are highlighted in bold. Key observations include Qwen2.5-72B surpassing GPT-4o in few-shot settings with an F1 score of 80.88\% at 10-shot, and the variability in optimal shot counts across models. Additionally, some models exhibit diminishing returns or inconsistent VCR at higher shot counts.}}
\resizebox{.7\textwidth}{!}{
\begin{tabular}{l|l|ccccccc}
\hline
Model & Metric & 0-shot & 5-shot & 10-shot & 20-shot & 30-shot & 40-shot & 50-shot \\
\hline
\multirow{2}{*}{GPT-4o-mini} & F1 & \textbf{72.96} & 71.81 & 69.17 & 67.98 & 69.13 & 69.25 & 71.05 \\
                             & VCR & 99.17 & \textbf{99.97} & 99.83 & 99.80 & 99.90 & 99.87 & 99.93 \\
\hline
\multirow{2}{*}{GPT-4o} & F1 & 79.53 & 80.00 & \textbf{80.36} & 79.67 & 80.31 & 79.93 & 80.14 \\
                        & VCR & 6.60 & 99.80 & \textbf{99.97} & 99.93 & 99.90 & 99.73 & 99.93 \\
\hline
\multirow{2}{*}{o1-mini} & F1 & 73.77 & 74.83 & 74.86 & 75.35 & 75.48 & 75.73 & \textbf{75.90} \\
                         & VCR & \textbf{99.90} & 99.67 & 99.43 & 99.47 & 99.77 & 99.77 & 99.77 \\
\hline
\multirow{2}{*}{o1-preview} & F1 & 74.51 & 75.35 & 76.77 & 76.25 & 76.44 & 76.74 & \textbf{77.08} \\
                            & VCR & \textbf{99.80} & 97.90 & 98.73 & 98.53 & 98.40 & 98.40 & 98.17 \\
\hline
\multirow{2}{*}{Llama3.1-8B} & F1 & 73.71 & 72.69 & 70.83 & 76.62 & \textbf{76.97} & 70.67 & 72.31 \\
                             & VCR & 80.33 & \textbf{98.87} & 96.23 & 97.90 & 73.27 & 75.90 & 66.23 \\
\hline
\multirow{2}{*}{Llama3.1-70B} & F1 & 75.26 & 75.45 & \textbf{75.67} & 73.49 & 75.65 & - & - \\
                              & VCR & 0.97 & 79.00 & \textbf{83.27} & 81.90 & 54.80 & - & - \\
\hline
\multirow{2}{*}{Mistral-7B} & F1 & 66.34 & \textbf{68.10} & 64.93 & 66.91 & 68.63 & 66.44 & - \\
                            & VCR & 1.17 & \textbf{14.20} & 10.63 & 8.27 & 7.00 & 6.33 & - \\
\hline
\multirow{2}{*}{InternLM2.5-7B} & F1 & 69.06 & \textbf{72.71} & 62.51 & 67.70 & 65.35 & 61.62 & 55.68 \\
                                & VCR & \textbf{100.00} & 99.90 & 98.83 & 94.73 & 94.03 & 98.13 & 98.03 \\
\hline
\multirow{2}{*}{InternLM2.5-7B-1M} & F1 & 52.45 & \textbf{77.01} & 66.58 & 67.64 & 68.10 & 71.25 & 70.31 \\
                                   & VCR & \textbf{99.87} & 94.53 & 88.67 & 82.13 & 82.37 & 83.37 & 88.47 \\
\hline
\multirow{2}{*}{InternLM2.5-20B} & F1 & \textbf{63.52} & - & - & - & - & - & - \\
                                 & VCR & \textbf{67.27} & - & - & - & - & - & - \\
\hline
\multirow{2}{*}{Qwen2-7B} & F1 & \textbf{71.01} & 43.65 & 59.91 & 60.67 & 67.83 & 69.25 & 65.34 \\
                              & VCR & \textbf{99.97} & 99.70 & 99.00 & 99.77 & 98.87 & 99.10 & 98.30 \\
\hline
\multirow{2}{*}{Qwen2-72B} & F1 & 75.72 & 77.22 & \textbf{77.40} & - & - & - & - \\
                               & VCR & 97.73 & 98.77 & \textbf{99.47} & - & - & - & - \\
\hline
\multirow{2}{*}{Qwen2.5-3B} & F1 & 55.06 & \textbf{64.16} & 64.03 & 52.65 & 53.11 & 62.15 & 64.52 \\
                            & VCR & \textbf{100.00} & 99.73 & 99.50 & 93.17 & 90.40 & 71.73 & 57.40 \\
\hline
\multirow{2}{*}{Qwen2.5-7B} & F1 & 61.01 & 62.19 & 67.91 & 72.78 & \textbf{75.28} & 74.72 & 74.14 \\
                            & VCR & \textbf{100.00} & 99.80 & 97.97 & 80.70 & 80.50 & 85.47 & 75.63 \\
\hline
\multirow{2}{*}{Qwen2.5-72B} & F1 & 76.99 & 80.40 & \textbf{80.88} & - & - & - & - \\
                             & VCR & \textbf{99.40} & 94.17 & 91.23 & - & - & - & - \\
\hline
\end{tabular}
}
\label{tab:performance}
\end{table*}

\subsection{Fine-tuning Open-source LLMs}
We fine-tuned open-source LLMs using LoRA for smaller models (3B, 7B, and 8B) and QLoRA for larger models (20B, 70B, and 72B). The results, presented in Table~\ref{tab:finetuning_performance}, reveal several important findings:

\subsubsection{Dramatic Enhancement in Valid Classification Ratio}
Fine-tuning significantly improved the Valid Classification Ratio (VCR) across all models. Remarkably, after just 0.2 epochs, nearly all models achieved and maintained a perfect VCR score of 100\%. This rapid improvement demonstrates that fine-tuning swiftly enhances the models' ability to produce valid classifications, effectively addressing a key limitation observed in few-shot prompting approaches.

\subsubsection{F1 Score Progression and Mitigating Overfitting}
Our initial experiments involved fine-tuning the LLMs for 6 epochs, with checkpoints saved at the conclusion of each epoch. We observed a common pattern: F1 scores for most models initially increased and reached a peak, but tended to decline after 1 epoch. This trend suggested potential overfitting, where models began to memorize training data at the expense of generalization.

To address this challenge, we refined our approach in subsequent experiments by implementing a more granular checkpoint strategy, saving model states every 0.2 epochs. This adjustment allowed us to:

\begin{itemize}
    \item Capture more detailed performance data throughout the training process.
    \item Identify and select the optimal results with greater precision.
    \item Effectively mitigate the overfitting observed in the later stages of training.
\end{itemize}

This refined approach yielded improved results. For instance, Qwen2-72B achieved its peak F1 score of 80.33\% at 0.8 epochs, demonstrating the benefits of our adjusted fine-tuning strategy in optimizing model performance while minimizing overfitting.

\subsubsection{Superior Performance of Fine-tuned Open-source LLMs}

Fine-tuning has enabled several open-source LLMs to achieve F1 scores that not only match but also surpass the best few-shot performance of GPT-4o, which previously stood at 80.36\%. Among the most notable results:

\begin{itemize}
    \item \textbf{Llama3.1-70B} surpassed GPT-4o with a robust F1 score of 80.77\% after 1.0 epoch.
    \item \textbf{Qwen2-72B} delivered an outstanding F1 score of 80.33\% after 0.8 epochs.
    \item \textbf{InternLM2.5-20B} and \textbf{InternLM2.5-7B-1M} both attained impressive F1 scores of 80.28\% after 0.8 epochs.
\end{itemize}

% These results highlight that fine-tuned open-source models can not only compete with but also outperform proprietary models in complex reasoning tasks.

% \subsubsection{Efficiency of Fine-tuning}
% The rapid improvement in both VCR and F1 scores within the first epoch highlights the efficiency of fine-tuning for enhancing model performance on specific tasks. This efficiency is particularly valuable given the computational constraints observed in few-shot evaluations of larger models.

% \subsubsection{Balancing Performance and Generalization}
% The observed pattern of increasing then decreasing F1 scores during fine-tuning emphasizes the importance of careful monitoring and potential early stopping to prevent overfitting. This underscores the need for robust fine-tuning strategies that can balance performance improvements with generalization capabilities.

These findings collectively demonstrate the power of fine-tuning in unlocking the full potential of open-source LLMs, enabling them to compete with and even surpass proprietary models in challenging reasoning tasks. 

\begin{table*}[htbp]
\centering
\small
\caption{Fine-tuning Performance of Open-Source Language Models Over Training Epochs (F1 Score and Valid Classification Ratio in \%). \small{The top metrics for each model across different epochs are highlighted in bold. Fine-tuning led to significant improvements, with most models achieving a perfect VCR of 100\% within the first 0.2 epochs. Llama3.1-70B outperformed GPT-4o with an F1 score of 80.77\%, followed closely by Qwen2-72B, InternLM2.5-7B-1M, and InternLM2.5-20B, all of which closely matched GPT-4o's performance.}}

\resizebox{0.9\textwidth}{!}{
\begin{tabular}{l|l|ccccccccccc}
\hline
Model & Metric & 0.0 & 0.2 & 0.4 & 0.6 & 0.8 & 1.0 & 1.2 & 1.4 & 1.6 & 1.8 & 2.0 \\
\hline
\multirow{2}{*}{Llama3.1-8B} & F1 & 73.71 & 74.28 & 74.89 & 73.71 & 78.64 & \textbf{79.25} & 76.29 & 77.45 & 76.62 & 75.80 & 76.17 \\
                             & VCR & 80.33 & \textbf{100.00} & 99.93 & \textbf{100.00} & \textbf{100.00} & \textbf{100.00} & \textbf{100.00} & \textbf{100.00} & \textbf{100.00} & \textbf{100.00} & \textbf{100.00} \\
\hline
\multirow{2}{*}{Llama3.1-70B} & F1 & 75.26 & 79.11 & 76.25 & 75.61 & 78.15 & \textbf{80.77} & 76.44 & 78.20 & 78.34 & 78.24 & 77.92 \\
                              & VCR & 0.97 & \textbf{100.00} & \textbf{100.00} & \textbf{100.00} & \textbf{100.00} & \textbf{100.00} & \textbf{100.00} & \textbf{100.00} & \textbf{100.00} & \textbf{100.00} & \textbf{100.00} \\
\hline
\multirow{2}{*}{Mistral-7B} & F1 & 66.34 & 73.43 & \textbf{75.37} & 70.68 & 74.05 & 74.72 & 73.49 & 76.48 & 74.10 & 72.57 & 73.08 \\
                            & VCR & 1.17 & \textbf{100.00} & \textbf{100.00} & \textbf{100.00} & \textbf{100.00} & \textbf{100.00} & \textbf{100.00} & \textbf{100.00} & \textbf{100.00} & \textbf{100.00} & \textbf{100.00} \\
\hline
\multirow{2}{*}{InternLM2.5-7B} & F1 & 69.06 & 73.30 & 73.97 & 71.04 & 76.60 & 75.14 & 74.89 & \textbf{76.53} & 74.32 & 73.47 & 74.08 \\
                                & VCR & \textbf{100.00} & \textbf{100.00} & \textbf{100.00} & \textbf{100.00} & \textbf{100.00} & \textbf{100.00} & \textbf{100.00} & \textbf{100.00} & \textbf{100.00} & \textbf{100.00} & \textbf{100.00} \\
\hline
\multirow{2}{*}{InternLM2.5-7B-1M} & F1 & 52.45 & 78.65 & 78.87 & 75.66 & \textbf{80.28} & 78.43 & 78.77 & 77.77 & 77.56 & 77.40 & 77.78 \\
                                   & VCR & 99.87 & \textbf{100.00} & \textbf{100.00} & \textbf{100.00} & \textbf{100.00} & \textbf{100.00} & \textbf{100.00} & \textbf{100.00} & \textbf{100.00} & \textbf{100.00} & \textbf{100.00} \\
\hline
\multirow{2}{*}{InternLM2.5-20B} & F1 & 63.52 & 76.93 & 78.49 & 74.22 & \textbf{80.28} & 79.38 & 77.87 & 79.65 & 79.00 & 77.95 & 78.27 \\
                                 & VCR & 67.27 & \textbf{100.00} & \textbf{100.00} & \textbf{100.00} & \textbf{100.00} & \textbf{100.00} & \textbf{100.00} & \textbf{100.00} & \textbf{100.00} & \textbf{100.00} & \textbf{100.00} \\
\hline
\multirow{2}{*}{Qwen2-7B} & F1 & 71.01 & 74.90 & 77.49 & 73.32 & 74.76 & \textbf{78.89} & 73.63 & 75.54 & 72.92 & 71.64 & 72.60 \\
                            & VCR & \textbf{100.00} & 99.97 & \textbf{100.00} & \textbf{100.00} & \textbf{100.00} & 99.93 & \textbf{100.00} & 99.97 & \textbf{100.00} & \textbf{100.00} & \textbf{100.00} \\
\hline
\multirow{2}{*}{Qwen2-72B} & F1 & 75.72 & 78.28 & 77.01 & 78.40 & \textbf{80.33} & 77.00 & 79.01 & 79.06 & 79.30 & 80.42 & 79.74 \\
                             & VCR & 97.73 & \textbf{100.00} & \textbf{100.00} & \textbf{100.00} & \textbf{100.00} & \textbf{100.00} & \textbf{100.00} & \textbf{100.00} & \textbf{100.00} & \textbf{100.00} & \textbf{100.00} \\
\hline
\multirow{2}{*}{Qwen2.5-3B} & F1 & 53.43 & 69.74 & 69.97 & 71.77 & 72.64 & 72.33 & 71.73 & \textbf{74.10} & 73.56 & 72.63 & 72.90 \\
                            & VCR & \textbf{100.00} & 99.90 & \textbf{100.00} & \textbf{100.00} & \textbf{100.00} & \textbf{100.00} & \textbf{100.00} & \textbf{100.00} & \textbf{100.00} & \textbf{100.00} & \textbf{100.00} \\
\hline
\multirow{2}{*}{Qwen2.5-7B} & F1 & 61.01 & 74.83 & 75.77 & 76.94 & 76.66 & \textbf{78.15} & 76.61 & 76.55 & 76.89 & 76.81 & 77.04 \\
                            & VCR & \textbf{100.00} & 99.80 & 99.97 & 99.97 & \textbf{100.00} & 99.93 & \textbf{100.00} & 99.93 & 99.97 & 99.93 & 99.97 \\
\hline
\multirow{2}{*}{Qwen2.5-72B} & F1 & 76.99 & 78.74 & 77.19 & 77.83 & \textbf{79.84} & 76.82 & 79.01 & 78.26 & 78.24 & 78.94 & 78.50 \\
                             & VCR & 99.40 & \textbf{100.00} & \textbf{100.00} & \textbf{100.00} & \textbf{100.00} & \textbf{100.00} & \textbf{100.00} & \textbf{100.00} & \textbf{100.00} & \textbf{100.00} & \textbf{100.00} \\
\hline
\end{tabular}
}
\label{tab:finetuning_performance}
\end{table*}

\subsection{Comparative Analysis of Proprietary and Open-source LLMs}

Table \ref{tab:comparing_performance} presents a comprehensive comparison of large language models (LLMs) across various evaluation settings: baseline (zero-shot), best few-shot, and best fine-tuned configurations. The "best" refers to the highest F1 score achieved in each category. The table is structured in two sections: the first lists OpenAI models without fine-tuned results, while the second features open-source models with fine-tuning data. This organization facilitates a clear comparison between proprietary and open-source models across different evaluation strategies.

\subsubsection{Fine-tuning vs. Few-shot Learning}

The results indicate that task-specific fine-tuning generally enhances model performance, especially for open-source models. Most open-source models demonstrated better performance after fine-tuning compared to their few-shot results. However, the Qwen2.5-72B model stands out as an exception, achieving the highest overall performance with an F1 score of 80.88\% via few-shot learning, surpassing GPT-4o's best few-shot score of 80.36\%. Llama3.1-70B secured the second-highest overall performance with an F1 score of 80.77\% through fine-tuning, also surpassing GPT-4o.

These findings suggest that while fine-tuning is generally advantageous, few-shot learning can sometimes lead to superior outcomes. Researchers and practitioners should carefully evaluate both strategies when optimizing model performance for specialized tasks.

\subsubsection{Extended Context Windows}

The performance of the InternLM2.5-7B-1M model underscores the significant potential of extended context windows, especially in models with relatively fewer parameters. Despite its modest baseline performance (F1 score of 52.45\%), the model showed remarkable improvement in advanced settings. In the few-shot configuration, InternLM2.5-7B-1M achieved an F1 score of 77.01\%, ranking third among all evaluated open-source LLMs. After fine-tuning, the model further improved to an F1 score of 80.28\%, making it the third-highest performer overall in the fine-tuned category.

These results suggest that extended context windows can significantly enhance model performance, particularly when combined with fine-tuning. This finding points to a promising direction for future model development, where leveraging extended context windows might lead to improved reasoning capabilities without necessitating a large increase in parameter count. Future research should explore how to best utilize extended context windows in model architectures to further advance LLMs' capacity for complex reasoning.


\begin{table*}[htbp]
\centering
\caption{Comparative Analysis of Large Language Models (LLMs) across Baseline (Zero-shot), Best Few-shot, and Best Fine-tuned Configurations (F1 Score and Valid Classification Ratio in \%). \small{"Best" refers to the highest F1 score achieved in each evaluation setting. The table is divided into two sections: OpenAI models (top) and open-source models (bottom). For each section, the highest metrics in each category are highlighted in bold. Notable observations include Qwen2.5-72B's superior few-shot performance with an F1 score of 80.88\%, surpassing GPT-4o's best score of 80.36\%. Among fine-tuned models, Llama3.1-70B achieves the highest F1 score of 80.77\%. The InternLM2.5-7B-1M model demonstrates significant improvement from a baseline F1 of 52.45\% to 77.01\% in few-shot and 80.28\% after fine-tuning, highlighting the potential of extended context windows.}}
\label{tab:comparing_performance}
\resizebox{0.6\textwidth}{!}{
\begin{tabular}{lcccccccccc}
\hline
\multirow{2}{*}{Model} & \multicolumn{2}{c}{Baseline (\%)} & \multicolumn{4}{c}{Best Few-shot (\%)} & \multicolumn{4}{c}{Best Fine-tuned (\%)} \\
\cmidrule(lr){2-3} \cmidrule(lr){4-7} \cmidrule(lr){8-11}
& F1 & VCR & F1 & VCR & Shots & & F1 & VCR & Epoch & \\
\hline
GPT-4o-mini & 72.96 & 99.17 & 72.96 & 99.17 & 0 & & - & - & - & \\
GPT-4o & \textbf{79.53} & 6.60 & \textbf{80.36} & \textbf{99.97} & 10 & & - & - & - & \\
o1-mini & 73.77 & \textbf{99.90} & 75.90 & 99.77 & 50 & & - & - & - & \\
o1-preview & 74.51 & 99.80 & 77.08 & 98.17 & 50 & & - & - & - & \\
\hline
Llama3.1-8B & 73.71 & 80.33 & 76.97 & 73.27 & 30 & & 79.25 & \textbf{100.00} & 1.0 & \\
Llama3.1-70B & 75.26 & 0.97 & 75.67 & 83.27 & 10 & & \textbf{80.77} & \textbf{100.00} & 1.0 & \\
Mistral-7B & 66.34 & 1.17 & 68.63 & 7.00 & 30 & & 76.48 & \textbf{100.00} & 1.4 & \\
InternLM2.5-7B & 69.06 & \textbf{100.00} & 72.71 & 99.90 & 5 & & 76.60 & \textbf{100.00} & 0.8 & \\
InternLM2.5-7B-1M & 52.45 & 99.87 & 77.01 & 94.53 & 5 & & 80.28 & \textbf{100.00} & 0.8 & \\
InternLM2.5-20B & 63.52 & 67.27 & - & - & - & & 80.28 & \textbf{100.00} & 0.8 & \\
Qwen2-7B & 71.01 & 99.97 & 71.01 & 99.97 & 0 & & 77.49 & \textbf{100.00} & 0.4 & \\
Qwen2-72B & 75.72 & 97.73 & 77.40 & 99.47 & 10 & & 80.33 & \textbf{100.00} & 0.8 & \\
Qwen2.5-3B & 55.06 & \textbf{100.00} & 64.16 & 99.73 & 5 & & 74.10 & \textbf{100.00} & 1.4 & \\
Qwen2.5-7B & 61.01 & \textbf{100.00} & 75.28 & 80.50 & 30 & & 78.15 & 99.93 & 1.0 & \\
Qwen2.5-72B & \textbf{76.99} & 99.40 & \textbf{80.88} & 91.23 & 10 & & 79.84 & \textbf{100.00} & 0.8 & \\
\hline
\end{tabular}
}
\end{table*}

